\cleardoublepage
\phantomsection
\begin{planbox}[แผนบริหารการสอนประจำบทที่ 11\quad กัมมันตภาพรังสีและการฉายรังสีควบคุมจุลินทรีย์ (สัปดาห์ที่ 12)]
\noindent\textbf{1. หัวข้อเนื้อหาประจำบท (บรรยาย 2 ชั่วโมง และปฏิบัติการ 3 ชั่วโมง)}
\begin{itemize}
    \item 11.1 โครงสร้างนิวเคลียส ไอโซโทปกัมมันตรังสี และชนิดของรังสี (แอลฟา บีตา แกมมา และรังสีอัลตราไวโอเลต)
    \item 11.2 แนวคิดครึ่งชีวิต ($T_{1/2}$) และการคำนวณจำนวนที่เหลืออยู่ด้วยการหารครึ่งทีละขั้นโดยไม่ใช้ลอการิทึม
    \item 11.3 การประยุกต์ใช้รังสีแกมมาและรังสี UV-C ($254\text{ nm}$) ในการทำลายดีเอ็นเอของจุลินทรีย์และกฎกำลังสองผกผันของความเข้มรังสี
    \item 11.4 ปฏิบัติการที่ 11 การจำลองการสลายตัวตามครึ่งชีวิตและการศึกษาผลของระยะห่างและเวลาฉายรังสี UV ต่อประสิทธิภาพการฆ่าเชื้อ
\end{itemize}

\noindent\textbf{2. วัตถุประสงค์เชิงพฤติกรรมและผลลัพธ์การเรียนรู้รายวิชา (CLOs ที่เชื่อมโยง PLOs)}
\begin{enumerate}
    \item เปรียบเทียบอำนาจทะลุทะลวงของรังสีแอลฟา บีตา แกมมา และรังสี UV-C และคำนวณครึ่งชีวิตด้วยวิธีหารสองเป็นขั้นๆ ได้ \textbf{(CLO1 เชื่อมโยง PLO2)}
    \item อธิบายเหตุผลทางฟิสิกส์ว่าทำไมรังสี UV-C จึงใช้ฆ่าเชื้อได้เฉพาะบนพื้นผิวและในตู้ปลอดเชื้อ (Biosafety Cabinet) แต่ไม่สามารถทะลุผ่านฝาพลาสติกของจานเพาะเชื้อได้ ในขณะที่รังสีแกมมาสามารถทะลุฆ่าเชื้อเวชภัณฑ์ทั้งกล่องได้ \textbf{(CLO1, CLO2 เชื่อมโยง PLO2, PLO3)}
    \item ปฏิบัติตามหลักความปลอดภัยทางรังสี 3 ประการ (เวลา ระยะห่าง และฉากกั้น) อย่างเคร่งครัด \textbf{(CLO3, CLO4 เชื่อมโยง PLO1, PLO5, PLO6)}
\end{enumerate}

\noindent\textbf{3. กิจกรรมการจัดการเรียนรู้เชิงรุก (Active Learning \& OBE)}
\begin{itemize}
    \item \textbf{Simulation \& Data Plotting}\quad ใช้การทอดลูกเต๋าหรือเหรียญจำลองการสลายตัวกัมมันตรังสีและการตายของประชากรแบคทีเรียเมื่อได้รับรังสี (Decimal Reduction Time, ค่า $D$-value) เพื่อสร้างกราฟการลดลงแบบครึ่งชีวิตด้วยตนเอง
\end{itemize}

\noindent\textbf{4. สื่อการเรียนรู้และแหล่งเรียนรู้}
\begin{itemize}
    \item เอกสารคำสอนบทที่ 11 ชุดจำลองครึ่งชีวิต และเครื่องวัดความเข้มแสง/รังสี UV ตามระยะห่าง
\end{itemize}

\noindent\textbf{5. การวัดและประเมินผลประจำบท}
\begin{itemize}
    \item แบบฝึกหัดการคำนวณครึ่งชีวิตและกฎกำลังสองผกผัน \textbf{(CLO1, CLO2)}
    \item กราฟครึ่งชีวิตและรายงานปฏิบัติการที่ 11 \textbf{(CLO3, CLO4)}
\end{itemize}
\end{planbox}
\clearpage

\chapter{กัมมันตภาพรังสีและการฉายรังสีควบคุมจุลินทรีย์}
\index{กัมมันตภาพรังสี}\index{ครึ่งชีวิต}\index{รังสีแกมมา}\index{รังสีอัลตราไวโอเลต UV-C}

\section{ชนิดของรังสีและอำนาจทะลุทะลวงในงานจุลชีววิทยา}

\deepskinfig{figures/sec_3d/ch11_s1.jpg}{กลไกการเกิดพันธะ Thymine Dimer ในสายดีเอ็นเอจากรังสี UV-C ($254\text{ nm}$) และการเปรียบเทียบอำนาจทะลุทะลวงของรังสีแอลฟา บีตา และแกมมา ($^{60}\text{Co}$)}

รังสีที่นำมาประยุกต์ใช้ในงานจุลชีววิทยาและการแพทย์แบ่งออกเป็น 2 กลุ่มใหญ่ตามพลังงาน ได้แก่ \textbf{รังสีไม่ก่อไอออน} (\textbf{Non-Ionizing Radiation}) เช่น \textbf{รังสีอัลตราไวโอเลตชนิดซี} (\textbf{UV-C} ความยาวคลื่นประมาณ $254\text{ nm}$) ซึ่งมีพลังงานพอเหมาะที่จะถูกดูดกลืนโดยสายดีเอ็นเอของแบคทีเรียและไวรัส ทำให้เกิดพันธะเชื่อมผิดปกติระหว่างเบสไทมีน (\textbf{Thymine Dimer}) จนจุลินทรีย์ไม่สามารถจำลองตัวเองและตายในที่สุด อย่างไรก็ตาม รังสี UV-C มีอำนาจทะลุทะลวงต่ำมาก แม้แต่แผ่นแก้วหรือฝาพลาสติกใสของจานเพาะเชื้อก็กั้นรังสี UV-C ได้เกือบหมด จึงใช้สำหรับการฆ่าเชื้อบนพื้นผิวโต๊ะปฏิบัติการและอากาศภายใน\textbf{ตู้ปลอดเชื้อชีวภาพ} (\textbf{Biosafety Cabinet}) เท่านั้น

ส่วนอีกกลุ่มหนึ่งคือ \textbf{รังสีก่อไอออน} (\textbf{Ionizing Radiation}) ที่เกิดจากการสลายตัวของนิวเคลียสกัมมันตรังสี ได้แก่ \textbf{รังสีแอลฟา} ($\alpha$ กั้นได้ด้วยแผ่นกระดาษ), \textbf{รังสีบีตา} ($\beta$ กั้นได้ด้วยแผ่นอะลูมิเนียมบาง) และ\textbf{รังสีแกมมา} ($\gamma$ เช่น จากโคบอลต์-60, $^{60}\text{Co}$) ซึ่งเป็นคลื่นแม่เหล็กไฟฟ้าพลังงานสูงมากที่มีอำนาจทะลุทะลวงผ่านกล่องพัสดุและพลาสติกได้อย่างสมบูรณ์ จึงนิยมใช้ฉายรังสีฆ่าเชื้อเข็มฉีดยา จานเพาะเชื้อพลาสติกสำเร็จรูป และอาหารส่งออก ดังแสดงใน\mbox{ภาพที่} 11.1

\begin{figure}[H]
\centering
\resizebox{0.88\textwidth}{!}{%
\begin{tikzpicture}[font=\sffamily, >=Stealth]
    % Sources on Left
    \foreach \y/\lbl/\col in {3.6/{รังสี UV-C (254 nm)}/labPurple, 2.5/{รังสีแอลฟา ($\alpha$)}/crimsonAlert, 1.4/{รังสีบีตา ($\beta$)}/physAmber, 0.3/{รังสีแกมมา ($\gamma$, $^{60}\text{Co}$)}/bioEmerald} {
        \fill[\col!20, draw=\col, line width=1.1pt, rounded corners=2pt] (0,\y-0.3) rectangle (3.5,\y+0.35);
        \node[font=\footnotesize\bfseries, color=\col!90!black] at (1.75,\y) {\lbl};
    }

    % Barrier 1: Plastic Petri Dish Lid / Paper
    \fill[physCyan!25, draw=physRoyalBlue, line width=1.2pt] (5.2,-0.2) rectangle (5.8,4.2);
    \node[above, font=\scriptsize\bfseries, align=center, color=physRoyalBlue] at (5.5,4.2) {ฝาพลาสติกจานเพาะเชื้อ\\หรือแผ่นกระดาษ};

    % Barrier 2: Aluminum / Thick Plastic Box
    \fill[physSlate!35, draw=physSlate, line width=1.2pt] (8.5,-0.2) rectangle (9.2,4.2);
    \node[above, font=\scriptsize\bfseries, align=center, color=charcoal] at (8.85,4.2) {แผ่นอะลูมิเนียม\\หรือกล่องเวชภัณฑ์};

    % Barrier 3: Thick Lead / Concrete
    \fill[charcoal!75, draw=charcoal, line width=1.4pt] (12.0,-0.2) rectangle (13.0,4.2);
    \node[above, font=\scriptsize\bfseries, align=center, color=charcoal] at (12.5,4.2) {กำแพงตะกั่วหนา\\หรือคอนกรีต};

    % Rays
    \draw[->, labPurple, line width=1.8pt, decorate, decoration={snake, amplitude=0.7mm, segment length=3mm}] (3.5,3.6) -- (5.2,3.6);
    \node[right, font=\scriptsize\bfseries, color=labPurple] at (5.9,3.6) {ถูกฝาพลาสติกกั้น (ฆ่าเชื้อเฉพาะผิวเปิดโล่ง)};

    \draw[->, crimsonAlert, line width=1.8pt] (3.5,2.5) -- (5.2,2.5);
    \draw[->, physAmber, line width=1.8pt] (3.5,1.4) -- (8.5,1.4);
    \draw[->, bioEmerald, line width=2.0pt, decorate, decoration={snake, amplitude=0.7mm, segment length=3mm}] (3.5,0.3) -- (12.0,0.3);
    \node[above, font=\scriptsize\bfseries, color=bioEmerald!90!black] at (10.5,0.4) {ทะลุฆ่าเชื้อจานเพาะเชื้อทั้งลัง};
\end{tikzpicture}%
}
\caption{การเปรียบเทียบอำนาจทะลุทะลวงของรังสี UV-C รังสีแอลฟา รังสีบีตา และรังสีแกมมาในการฆ่าเชื้อทางจุลชีววิทยา}
\end{figure}

\vspace{4pt}
\arucomarkerbadge{10}{11}{กัมมันตภาพรังสี กฎกำลังสองผกผัน และการฆ่าเชื้อด้วยรังสี UV-C (254 nm)}{จำลองการแผ่รังสี UV-C ทำลายพันธะไทมีนในสายดีเอ็นเอ (Thymine Dimer) การลดลงของความเข้มรังสีตามระยะทาง $I \propto 1/r^2$ และครึ่งชีวิตกัมมันตภาพรังสี}{ระยะห่างจากหลอดรังสี $r$ ($10\text{--}100\text{ cm}$), เวลาฉายรังสี $t$ ($\text{s}$), อัตราการรอดชีวิตของจุลินทรีย์ (\%)}


\section{การคำนวณครึ่งชีวิตแบบเข้าใจง่ายโดยไม่ต้องใช้ลอการิทึม}

\deepskinfig{figures/sec_3d/ch11_s2.jpg}{การลดลงของปริมาณไอโซโทปกัมมันตรังสีและอัตราการรอดชีวิตของจุลินทรีย์แบบหารครึ่งทีละรอบครึ่งชีวิต ($T_{1/2}$) และกฎกำลังสองผกผันของความเข้มรังสี ($I \propto 1/r^2$)}

\textbf{ครึ่งชีวิต} (\textbf{Half-Life}, $T_{1/2}$) คือระยะเวลาที่สารกัมมันตรังสีสลายตัวจนเหลือปริมาณเพียง\textbf{ครึ่งหนึ่ง} ($50\%$ หรือ $\frac{1}{2}$) ของปริมาณเริ่มต้น นักศึกษาไม่จำเป็นต้องใช้ฟังก์ชันเอกซ์โพเนนเชียลหรือลอการิทึมในการคำนวณ แต่สามารถใช้\textbf{วิธีนับจำนวนรอบของการหารครึ่ง} ($n = \text{เวลาทั้งหมด } t \div T_{1/2}$) ได้ทันทีตามสมการที่ (11.1)

\begin{equation}
    N = N_0 \times \left(\frac{1}{2}\right)^n \qquad \text{โดยที่} \qquad n = \frac{t}{T_{1/2}}
\end{equation}

กล่าวคือ เมื่อผ่านไป $1$ ครึ่งชีวิตจะเหลือ $\frac{1}{2}$ ($50\%$), ผ่านไป $2$ ครึ่งชีวิตจะเหลือ $\frac{1}{4}$ ($25\%$), ผ่านไป $3$ ครึ่งชีวิตจะเหลือ $\frac{1}{8}$ ($12.5\%$) และผ่านไป $4$ ครึ่งชีวิตจะเหลือ $\frac{1}{16}$ ($6.25\%$) ตามลำดับ

นอกจากนี้ ในการใช้งานหลอดรังสี UV-C เพื่อฆ่าเชื้อ ความเข้มของรังสี ($I$) ที่ตกกระทบบนพื้นผิวจะลดลงตาม\textbf{กฎกำลังสองผกผันของระยะห่าง} (\textbf{Inverse Square Law}) คือ $I \propto 1/r^2$ ดังนั้น หากวางจานเพาะเชื้อห่างจากหลอด UV เพิ่มขึ้นเป็น $2$ เท่า ความเข้มรังสี UV ที่ตกกระทบจะลดลงเหลือเพียง $1/2^2 = 1/4$ เท่าของเดิม

\begin{worked}[การคำนวณครึ่งชีวิตของไอโซโทปติดตามทางชีววิทยาด้วยวิธีหารสองทีละขั้น]
\textbf{โจทย์ปัญหา}\quad นักวิจัยจุลชีววิทยาใช้ไอโซโทปฟอสฟอรัส-32 ($^{32}\text{P}$) ซึ่งมีครึ่งชีวิต $T_{1/2} = 14\text{ วัน}$ ในการติดฉลากสายดีเอ็นเอของแบคทีเรีย หากเริ่มต้นมีสารกัมมันตรังสี $^{32}\text{P}$ มวล $N_0 = 80\text{ }\mu\text{g}$ เมื่อเก็บไว้เป็นเวลา $t = 42\text{ วัน}$ จะเหลือ $^{32}\text{P}$ อยู่กี่ไมโครกรัม ($\mu\text{g}$)

\textbf{วิธีทำทีละขั้นตอนฉบับเข้าใจง่าย}
\begin{enumerate}
    \item \textbf{ขั้นที่ 1 หาจำนวนรอบของครึ่งชีวิต ($n = t / T_{1/2}$)}
    \[
        n = \frac{42\text{ วัน}}{14\text{ วันต่อรอบ}} = 3\text{ รอบครึ่งชีวิต}
    \]
    \item \textbf{ขั้นที่ 2 นำค่าเริ่มต้น $80\text{ }\mu\text{g}$ มาหารด้วย $2$ จำนวน $3$ ครั้งต่อเนื่องกัน}
    \[
        80\text{ }\mu\text{g} \xrightarrow{\text{รอบที่ 1 (14 วัน)}} 40\text{ }\mu\text{g} \xrightarrow{\text{รอบที่ 2 (28 วัน)}} 20\text{ }\mu\text{g} \xrightarrow{\text{รอบที่ 3 (42 วัน)}} 10\text{ }\mu\text{g}
    \]
\end{enumerate}
\textbf{คำตอบ}\quad เมื่อเวลาผ่านไป $42\text{ วัน}$ จะเหลือไอโซโทป $^{32}\text{P}$ อยู่เท่ากับ $10\text{ }\mu\text{g}$
\end{worked}

\section{ปฏิบัติการประจำสัปดาห์ที่ 12}

\deepskinfig{figures/sec_3d/ch11_s3.jpg}{ชุดปฏิบัติการที่ 11 สำหรับจำลองกราฟการสลายตัวครึ่งชีวิตและวัดความเข้มแสงตามระยะห่างเพื่อตรวจสอบกฎกำลังสองผกผัน}

\begin{labbox}[ปฏิบัติการที่ 11\quad การจำลองครึ่งชีวิตกัมมันตภาพรังสีและประสิทธิภาพรังสี UV (3 ชั่วโมง)]
\noindent\textbf{ผลลัพธ์การเรียนรู้ที่ประเมิน}\quad CLO1 (ครึ่งชีวิตและกฎกำลังสองผกผัน), CLO2 (การควบคุมจุลินทรีย์ด้วยรังสี UV), CLO3 (ความปลอดภัยทางรังสี), CLO4 (การเขียนกราฟครึ่งชีวิต)

\vspace{4pt}
\noindent\textbf{1. วัสดุและอุปกรณ์}
\begin{itemize}
    \item เหรียญหรือลูกเต๋าจำลองอะตอมกัมมันตรังสีและเซลล์แบคทีเรียจำนวน 100 ชิ้น พร้อมกล่องเขย่า
    \item แหล่งกำเนิดแสงจำลองหลอด UV และเครื่องวัดความเข้มแสงดิจิทัล (Lux / UV Radiometer) พร้อมไม้บรรทัดวัดระยะห่าง
\end{itemize}

\noindent\textbf{2. ขั้นตอนการทดลอง}
\begin{enumerate}
    \item \textbf{ตอนที่ 1 การจำลองกราฟการสลายตัวแบบครึ่งชีวิต}\quad เริ่มต้นด้วยเหรียญจำนวน $N_0 = 100$ เหรียญ เขย่าในกล่องแล้วเทออก คัดเหรียญที่หงายหน้าก้อยออก (ถือว่าสลายตัวหรือเซลล์ตายจากรังสี) นับจำนวนเหรียญที่เหลือรอด ($N$) บันทึกผลแล้วนำเหรียญที่เหลือมาเขย่าซ้ำอีก 5 รอบ เขียนกราฟระหว่างจำนวนที่เหลือ $N$ กับจำนวนรอบการเขย่า
    \item \textbf{ตอนที่ 2 การตรวจสอบกฎกำลังสองผกผันของความเข้มรังสีตามระยะห่าง}\quad วัดค่าความเข้มแสงเมื่อวางหัววัดห่างจากหลอดกำเนิดแสงเป็นระยะ $r = 10\text{ cm}, 20\text{ cm}$ และ $40\text{ cm}$ เปรียบเทียบอัตราส่วนความเข้มที่ลดลง
\end{enumerate}

\noindent\textbf{3. ตารางบันทึกผลการทดลอง}
\begin{center}
\footnotesize\setlength{\tabcolsep}{3.5pt}
\begin{tabularx}{\linewidth}{c c c c c c c}
\toprule
\textbf{รอบที่ ($n$)} & \textbf{เริ่มต้น ($n=0$)} & \textbf{รอบที่ 1} & \textbf{รอบที่ 2} & \textbf{รอบที่ 3} & \textbf{รอบที่ 4} & \textbf{รอบที่ 5} \\
\midrule
จำนวนตามทฤษฎี & $100$ & $50$ & $25$ & $12.5$ & $6.25$ & $3.125$ \\
จำนวนที่วัดได้จริง & $100$ & ........ & ........ & ........ & ........ & ........ \\
\bottomrule
\end{tabularx}
\end{center}
\end{labbox}

\section*{คำถามท้ายบทที่ 11 (ทบทวนแนวคิดและการคำนวณเบื้องต้น)}
\begin{enumerate}
    \item เพราะเหตุใดขณะเปิดหลอดรังสี UV-C ในตู้ปลอดเชื้อชีวภาพ (Biosafety Cabinet) นักศึกษาจึงต้องปิดกระจกหน้าตู้ลงและห้ามสอดมือหรือจ้องมองแสงสีม่วงของหลอด UV โดยตรง
    \item หากการฉายรังสี UV-C ทุกๆ $2\text{ นาที}$ สามารถลดจำนวนแบคทีเรียบนพื้นผิวลงได้ครึ่งหนึ่ง ($50\%$) หากเริ่มต้นมีแบคทีเรีย $1{,}600\text{ เซลล์}$ เมื่อเปิดไฟฉายรังสีนาน $8\text{ นาที}$ จะเหลือแบคทีเรียรอดชีวิตกี่เซลล์
\end{enumerate}
